\documentclass[10pt,a4paper]{amsart}
\usepackage[margin=19mm]{geometry}
\usepackage{amsmath,amssymb,mathtools}
\usepackage[scr=rsfs,cal=euler]{mathalfa}
\usepackage{xfrac}
\usepackage[unicode,hidelinks,bookmarks=false]{hyperref}
\newcommand{\ie}{i.e.\ }
\newcommand{\cf}{cf.\ }
\newcommand{\resp}{respectively\ }
\newcommand{\Subsection}{\subsection}
\providecommand{\nobreakdash}{\nobreak\hbox{-}}
\setlength{\emergencystretch}{2em}
\multlinegap0pt
\usepackage{bm}
\usepackage{bbm}
\usepackage{dsfont}
\usepackage{xspace}
\usepackage{cancel}
\usepackage{mathtools}
\usepackage{amsthm}
\makeatletter
\@ifundefined{theorem}{\newtheorem{theorem}{Theorem}}{}
\makeatother
\newcommand{\bx}{\boldsymbol{x}}
\title[Persistence, Thresholds, and Trait Composition in a Regulate]{Persistence, Thresholds, and Trait Composition in a Regulated Mutation-Selection Model}
\author{Phil. Pollett}
\date{Source: arXiv:2606.31146v1 (CC BY 4.0)}
\begin{document}
\maketitle

\noindent\emph{Source and licence.} Persistence, Thresholds, and Trait Composition in a Regulated Mutation-Selection Model. Version arXiv:2606.31146v1 (CC BY 4.0), distributed under the Creative Commons Attribution 4.0 International licence, as recorded in the arXiv licence metadata for this exact version and shown on the abstract page https://arxiv.org/abs/2606.31146v1. Full rights notice: licenses/arxiv-2606.31146-CC-BY-4.0.txt.\par\smallskip

\noindent\textit{Editorial scope.} Counted unit: the Theorem whose source label is \texttt{None} in the article (source0.tex lines 1581--1583, source label \texttt{None}), delivered together with the article's own complete original proof of it (source0.tex lines 1585--1644, one \texttt{proof} environment of 60 lines). No mathematics is written, repaired or re-derived: statement and proof are the article's own text line for line. Required context retained uncounted: none. Every definition, display and label of those lines is retained unchanged. Omitted from this package and not redistributed: 1--1580, 1584--1584, 1645--2028. Nothing inside a retained excerpt is omitted or altered: every retained body is delivered exactly as the article's lines are written, so no omission receipt of any kind accompanies this package. Citations: the retained excerpts carry no citation command at all, so no bibliography entry of the article is transcribed and no reference is redistributed. Reference adaptation: none was needed and none was made; every reference inside the delivered unit resolves inside it, so no unresolved reference or citation is produced. Printed slips kept literal: no printed word, symbol, hypothesis or number of the counted statement or of its proof was altered. Layout and class: the article's own class is \texttt{article} with packages including amsmath, amssymb, mathtools, amsthm, natbib, graphicx, epstopdf; the wrapper is the lane's pinned \texttt{amsart} editorial wrapper at 10pt on A4, which restates the theorem-like environments the retained text needs and adds the article's own macro definitions that the retained text needs (\texttt{\textbackslash bx}), with extra wrapper packages (bm, bbm, dsfont, xspace, cancel, mathtools). The delivered numbering is the unit's own. Assets: no retained span contains a figure environment, a graphics-inclusion command, a PDF-page inclusion, a table or a TikZ picture; no image file is copied into this package, the package declares no asset dependency and no image file is redistributed with it. Licence: version arXiv:2606.31146v1 of this article is distributed under the Creative Commons Attribution 4.0 International licence, as recorded in the arXiv licence metadata for this exact version and shown on the abstract page https://arxiv.org/abs/2606.31146v1. A full rights notice is delivered at licenses/arxiv-2606.31146-CC-BY-4.0.txt. Characters: every retained excerpt is pure ASCII, so the delivered engine, XeLaTeX with its default Latin Modern text font, reports no missing character.
\begin{theorem}
If $\bx_0\in E$, then $\bx(t)\in E$ for all $t>0$.
\end{theorem}

\begin{proof}
Integrating $\dot{\bx}=G(\bx) - D\bx$ using the integrating factor
$e^{Dt}$ we get
$$
\bx(t)=e^{-D t}\bx(0)+\int_0^t e^{-D(t-s)}\,G(\bx(s))\,ds,
$$
or, componentwise,
\begin{equation}
x_i(t)=e^{-\mu_i t}x_i(0)+\int_0^t
e^{-\mu_i(t-s)}\,\lambda\bigl(1-m(\bx(s))\bigr)\, (P\bx(s))_i\,ds,
\qquad i\in \{0,1\}.
\label{mildsolution1}
\end{equation}

Suppose that $\bx(0)\in E$, that is, 
$x_i(0)\geq 0$ for both $i$, and $m(\bx(0))\leq 1$.
We will first show $m(\bx(t))\leq 1$ for all $t$.
Let $t_*:=\inf\{t>0:\ m(\bx(t))>1\}$ 
($t_*=\infty$ if the set is empty).
Fix $t<t_*$. Then $m(\bx(s))\le 1$ for all $s\in[0,t]$, and hence 
$1-m(\bx(s))\geq 0$ on this interval. 
On summing (\ref{mildsolution1}) over $i$ we find that 
\[
m(\bx(t))=\sum_i e^{-\mu_i t}x_i(0)
+\int_0^t \lambda(1-m(\bx(s)))\sum_i e^{-\mu_i(t-s)}(P\bx(s))_i\,ds.
\]
But, $\sum_i e^{-\mu_i t}x_i(0)\le \sum_i x_i(0)=m(\bx(0))$, and
$$
\sum_i e^{-\mu_i(t-s)}(P\bx(s))_i
=\sum_k x_k(s)\sum_i e^{-\mu_i(t-s)}p_{i-k}\\
\leq \sum_k x_k(s)=m(\bx(s)),
$$
and hence 
\begin{equation}
m(\bx(t))\le m(\bx(0))+\int_0^t \lambda(1-m(\bx(s)))\,m(\bx(s))\,ds.
\label{mbxt1}
\end{equation}
Now compare the right-hand side of~(\ref{mbxt1}) with the
integrated form of the solution to the
logistic equation $\dot{u}=f(u)$, where $f(u)=\lambda u(1-u)$ and 
$u(0)=u_0:=m(\bx(0))\le 1$:
$$
u(t) = u_0 + \int_0^t \lambda (1-u(s)) u(s)\, ds.
$$
We will show that $m(\bx(t))\leq u(t)$ for all $t\geq 0$. 
Since $u(t)<1$ for all~$t$, it will follow that $t_*=\infty$. 
Let $\phi_t(x_0)$ be the flow corresponding to $\dot{u}=f(u)$,
and notice that the flow is monotone in that if $u_2\geq u_1>0$,
then $\phi_t(u_2)\geq \phi_t(u_1)$ for all~$t$.
Let $\tau := \inf\{t>0 : m(\bx(t)) > u(t)\}$.
Then, $m(\bx(t))\leq u(t)$ for all $t<\tau$, and $m(\bx(\tau))=u(\tau)$
by continuity. Now fix $h>0$ and observe that
$$
m(\bx(\tau+h)) \leq \phi_h(m(\bx(\tau)))=\phi_h(u(\tau))= u(\tau+h),
$$
which is a contradiction. Therefore $m(\bx(t))\leq u(t)$.
It is now easy to see from (\ref{mildsolution1}) that $x_i(t)\geq 0$
for all $i$, since the integrand is non-negative.
So, finally, $\bx(t)\in E$ for all $t>0$.
\end{proof}

\end{document}
